\subsection{Quantitative bounds from positive Hilbert data}

\begin{proposition}[P-F9: a common-period bound for the positivity index]
\label{prop:periodbound}
Let $f(s)$, $s\ge0$, be a quasi-polynomial whose generating function has
an explicit positive form
\[
 \sum_{s\ge0}f(s)T^s=\frac{H(T)}{\prod_{j=1}^r(1-T^{a_j})},
 \quad H(T)\in\mathbb Z_{\ge0}[T]\setminus\{0\},\quad a_j>0.
\]
Put $h=\deg H$, let $L$ be any positive common multiple of the $a_j$, and suppose
$r\ge2$. In the definitions of GCT VI, Definitions 1.2.2 and 1.2.4,
\[
 s(f)\le p(f)\le
 \max\{0,h+(r-1)L-\textstyle\sum_j a_j\}
 \le (r-1)L-1.
 \tag{P-F9}
\]
For $r=1$, $f$ is periodic with nonnegative constant constituents and
$p(f)=s(f)=0$. Thus a positive form with polynomially bounded $r$ and
common multiple $L$ supplies polynomial positivity and saturation
indices. A bound on the largest $a_j$ alone does not make this particular
estimate polynomial, because their least common multiple can be large.
\end{proposition}
\begin{proof}
First $h<\sum_j a_j$. To see this without an assumption on cancellations,
write the quasi-polynomial on each residue class modulo a period as a
polynomial in its nonnegative quotient variable. Express that polynomial
in the binomial basis $\binom{k}{j}$. Its generating function is a linear
combination of $T^{c+jD}/(1-T^D)^{j+1}$ with $0\le c<D$, all proper
rational functions. Their sum tends to zero at infinity. The displayed
positive rational expression for the same function therefore also has
numerator degree strictly below its denominator degree.

Multiply numerator and denominator by geometric sums to obtain
\[
 \frac{B(T)}{(1-T^L)^r},\qquad
 B(T)=H(T)\prod_{j=1}^r
       (1+T^{a_j}+\cdots+T^{L-a_j}).
\]
All coefficients $b_k$ of $B$ are nonnegative and
$D_B:=\deg B=h+rL-\sum_j a_j\le rL-1$.
For a fixed residue $c$ modulo $L$, the eventual constituent of $f$ is
\[
 F_c(s)=\sum_{k\equiv c\ (L)} b_k
 \binom{(s-k)/L+r-1}{r-1}
 =\frac1{L^{r-1}(r-1)!}
 \sum_{k\equiv c\ (L)} b_k
 \prod_{j=1}^{r-1}(s-k+jL).
\]
Indeed this follows by expanding $(1-T^L)^{-r}$ for all sufficiently
large $s$. The constituents coincide with those of the given
quasi-polynomial everywhere: two polynomials agreeing on an infinite
arithmetic progression agree identically (use a common period if
necessary). If $p\ge\max\{0,D_B-L\}$, each factor of $F_c(s+p)$ is
$s+(p-k+jL)$ with nonnegative constant term. Every constituent thus has
nonnegative coefficients. Translation merely permutes the residue
labels, so this proves the bound on $p(f)$. Any nonzero polynomial with
nonnegative coefficients is positive at every real $s\ge1$, which gives
$s(f)\le p(f)$. For $r=1$ the constituents are nonnegative constants.
Finally $D_B\le rL-1$ gives the last bound when $r\ge2$.
\end{proof}

This is an explicit sufficient condition for the revised PH2, not a proof
of that hypothesis for the GCT families. Taking $L$ to be the least common
multiple always gives a bound whose bit length is polynomial in the
expanded positive-form data. GCT VI, Lemma 3.1.6, supplies a different
polynomial modular-index bound for saturation in its stated setting;
we do not replace that known result by a new claim. The remaining
quantitative issue is to control the appropriate parameters for the
specific multiplicity rings, uniformly in the GCT input.

\begin{proposition}[P-F10: purity and normality do not control a grading's indices]
\label{prop:weightedplane}
Let $a,b>1$ be relatively prime. Give the polynomial ring $k[x,y]$ the
positive grading $\deg x=a$, $\deg y=b$, and put
$f(n)=\dim_k k[x,y]_n$. Then $f$ is a quasi-polynomial and its GCT
positivity and saturation indices are exactly
\[
 p(f)=s(f)=ab-a-b.
 \tag{P-F10}
\]
The underlying variety is always the same smooth normal affine plane.
Over a finite field its only compactly supported cohomology is
$H_c^4=\mathbb Q_\ell(-2)$; its intersection complex is the shifted
constant sheaf. Thus even the same pure geometry can support arbitrarily
large indices when its grading varies.
\end{proposition}
\begin{proof}
The monomial basis gives
$f(n)=\#\{(u,v)\in\mathbb Z_{\ge0}^2:au+bv=n\}$ and generating function
$1/((1-T^a)(1-T^b))$. Write
\[
 \frac1{(1-T^a)(1-T^b)}
 =\frac{\displaystyle\sum_{i=0}^{b-1}\sum_{j=0}^{a-1}T^{ai+bj}}
 {(1-T^{ab})^2}.
\]
The $ab$ exponents $r=ai+bj$ represent distinct residues modulo $ab$:
reduction modulo $b$ fixes $i$, and modulo $a$ fixes $j$.
Their minimum is zero and their maximum is
$r_{\max}=a(b-1)+b(a-1)=2ab-a-b<2ab$.
For the residue represented by $r$, its constituent is therefore
\[
 F_r(n)=\frac{n-r}{ab}+1.
\]
For $n\ge r$ this is the geometric-series coefficient.
If $0\le n<r$ in the same residue, necessarily $n=r-ab$;
the displayed formula is zero, as the coefficient is. Hence the
constituent formula is valid for every nonnegative integer $n$.

All constituents have positive slope $1/(ab)$. After a shift by $p$,
the smallest constant coefficient is
$1+(p-r_{\max})/(ab)=(p-(ab-a-b))/(ab)$.
It is nonnegative precisely when $p\ge ab-a-b$, proving the positivity
index. Every constituent is nonzero. Strict saturation after a shift
$s$ requires its value at $n=1$ to be strictly positive for the worst
constituent; the condition is $1+s>ab-a-b$, or, for integer $s$,
$s\ge ab-a-b$. This proves the second index.

The ring is a polynomial ring independently of its grading, hence its
spectrum is smooth and normal. The cohomology statement is the standard
affine-space calculation and duality (Milne, \emph{Lectures on \'Etale
Cohomology}, Example 16.3 and Theorem 24.1). Smoothness identifies its
perverse intersection complex with $\mathbb Q_\ell[2]$.
Taking $a=2$ and odd $b$ arbitrarily large makes the indices $b-2$
unbounded while this underlying geometry stays fixed.
\end{proof}

P-F10 by itself makes no representation-theoretic identification.
P-F15 below supplies that identification for the broad subgroup
restriction problem, where it refutes the printed rank-only index bounds.
It does not refute the narrower irreducible-inner plethysm bounds.
Independently, it refutes a proposed derivation of quantitative bounds
from purity, smoothness or normality alone. The action and its
grading, and a bound relating their weights to the input, are essential
additional data. Normality of the original ring also must not be
confused with saturation of the image of its exponent semigroup under
the degree map.
